\clearpage

\thispagestyle{empty}
\DEoEN{%
{\LARGE\textsf{\textbf{Erklärung zur Verwendung generativer KI-Systeme}}\bigskip}

Bei der Erstellung der eingereichten Arbeit habe ich auf künstlicher Intelligenz (KI) basierte Systeme benutzt: 

\begin{enumerate}
\item[$\boxtimes$] ja
\item[$\square$] nein\footnote{%
Die Erklärung ist in jedem Fall zu unterzeichnen, auch wenn Sie keine KI-Systeme genutzt haben und Ihr Kreuz bei \enquote{nein} gesetzt haben.}
\end{enumerate}

% Falls nein: Rest des Blocks entfernen, Unterschrift aber beibehalten

Falls ja: Die nachfolgend aufgeführten auf künstlicher Intelligenz (KI) basierten Systeme habe ich bei der Erstellung der eingereichten Arbeit benutzt:

\begin{enumerate}
\item Claude (Anthropic)
\item ChatGPT (OpenAI)
\item Microsoft Copilot
\end{enumerate}

Ich erkläre, dass ich

\begin{itemize}
  \item mich aktiv über die Leistungsfähigkeit und Beschränkungen der oben genannten 
KI-Systeme informiert habe,\footnote{U.a. gilt es hierbei zu beachten, dass an KI weitergegebene Inhalte ggf. als Trainingsdaten genutzt und wiederverwendet werden. Dies ist insb. für betriebliche Aspekte als kritisch einzustufen.}
  \item die aus den oben angegebenen KI-Systemen direkt oder sinngemäß übernommenen Passagen gekennzeichnet habe,
%
% In der Fußnote Ihrer Arbeit geben Sie die KI als Quelle an, z.B.: 
% Erzeugt durch Microsoft Copilot am dd.mm.yyyy. 
% Oder: Entnommen aus einem Dialog mit Perplexity vom dd.mm.yyyy. 
% Oder: Absatz 2.3 wurde durch ChatGPT sprachlich geglättet.
%
  \item überprüft habe, dass die mithilfe der oben genannten KI-Systeme generierten und von mir übernommenen Inhalte faktisch richtig sind,
  \item mir bewusst bin, dass ich als Autorin bzw. Autor dieser Arbeit die Verantwortung für die in ihr gemachten Angaben und Aussagen trage.
\end{itemize}

\clearpage

Die oben genannten KI-Systeme habe ich wie im Folgenden dargestellt eingesetzt: 

\begin{tabular}{|p{3.5cm}|p{4.3cm}|p{6.2cm}|}
\hline
\textbf{Arbeitsschritt} & \textbf{KI-System(e)} & \textbf{Verwendungsweise} \\
\hline

Ideenfindung und Strukturierung &
Claude (Anthropic) &
Unterstützung bei der Strukturierung der Arbeit, Formulierung von Forschungsfragen sowie Gliederung einzelner Kapitel. Inhalte wurden kritisch geprüft und eigenständig überarbeitet. \\
\hline

Sprachliche Überarbeitung und Glättung &
ChatGPT (OpenAI) &
Unterstützung bei der sprachlichen Verbesserung einzelner Textpassagen (Grammatik, Stil). Inhaltliche Aussagen wurden nicht automatisiert übernommen. \\
\hline

Code-Unterstützung &
GitHub Copilot (Microsoft) &
Unterstützung bei der Erklärung und Optimierung von Codefragmenten im Rahmen der Implementierung. Der finale Code wurde eigenständig entwickelt und validiert. \\
\hline

\end{tabular}
} % Ende deutscher Teil
{% Beginn englische Erklaerung
{\LARGE\textsf{\textbf{Declaration on the Use of Generative AI Systems}}\bigskip}

In preparing the submitted work, I have used artificial intelligence (AI)-based systems:

\begin{enumerate}
\item[$\boxtimes$] yes
\item[$\square$] no\footnote{%
The declaration must be signed in any case, even if you have not used any AI systems and have checked the box \enquote{no}.}
\end{enumerate}

% Falls nein: Rest des Blocks entfernen, Unterschrift aber beibehalten

If yes: I have used the following AI-based systems:


\begin{enumerate}
\item
\item
\item \ldots
\end{enumerate}

I hereby declare that I

\begin{itemize}
  \item actively informed myself about the capabilities and limitations of the above-mentioned AI systems,\footnote{In particular, it should be noted that content passed on to AI may be used as training data and reused. This is to be considered critical, especially for operational aspects.}
  \item indicated passages directly or indirectly adopted from the above-mentioned AI systems,
%
% In the footnote of your work, cite the AI as the source, e.g.:
% Generated by Microsoft Copilot on mm.dd.yyyy. 
% Or: Taken from a dialogue with Perplexity on mm.dd.yyyy. 
% Or: Paragraph 2.3 was linguistically smoothed by ChatGPT.
%
  \item verified that the content generated and adopted by me using the above-mentioned AI systems is factually correct,
  \item am aware that as the author of this work, I bear responsibility for the statements and information provided in it.
\end{itemize}

I have utilized the above-mentioned AI systems as illustrated below: 

\begin{center}
\begin{tabular}{|p{4cm}|p{3cm}|p{7cm}|}
    \hline
    \textbf{Task in the scientific work} &
%
% Examples of such tasks include: idea generation, 
% conception of the work, literature search, literature analysis, 
% literature management, selection of methods, data collection, 
% data analysis, generation of program code
%
% If you are unsure whether you need to specify an AI system used, 
% consult your supervisor.
%
	 \textbf{AI System(s) Used} & \textbf{Description of Usage} \\
    \hline
    & & \\ % Placeholder for your entries, insert your information here
    \hline
    & & \\ % (Expand the table as needed and continue on subsequent pages)
    \hline
    & & \\
    \hline
    & & \\
    \hline
  \end{tabular}
\end{center}
}

\vspace{3cm}

\begin{center}
\begin{tabular}{ccc}
(\DEoEN{Ort, Datum}{place, date}) & \hspace{0.3\linewidth} & (\DEoEN{Unterschrift}{signature})
\end{tabular}
\end{center}